\documentclass[12pt]{article}
\usepackage[utf8]{inputenc}
\usepackage[T1]{fontenc}
\usepackage[margin=1in]{geometry}
\usepackage{setspace}
\usepackage{xcolor}
\usepackage[colorlinks=true, urlcolor=blue!60!black, linkcolor=blue!60!black]{hyperref}
\setlength{\parindent}{2em}
\setlength{\parskip}{1.25em}
\renewcommand{\baselinestretch}{1.0}
\title{\textit{R.A.V. v. City of St. Paul}, 505 U.S. 377 (1992)}
\author{Case Brief}
\date{}
\begin{document}
\maketitle

% R.A.V. v. City of St. Paul, 505 U.S. 377 (1992)
% Chapter: Free Speech & Press

\section{Case Metadata}
\begin{description}
  \item[Case Name \& Citation:] \textit{R.A.V. v. City of St. Paul}, 505~U.S.~377 (1992) (No.~90--7675).
  \item[Argued / Decided:] Argued December~4, 1991; Decided June~22, 1992.
  \item[Vote:] 9--0 judgment (reversing and remanding \textit{In re Welfare of R.A.V.}, 464~N.W.2d~507 (Minn.~1991)); 5--4 on rationale.
  \item[Author:] Justice Scalia (joined by Chief Justice Rehnquist and Justices Kennedy, Souter, and Thomas).
  \item[Concurrences / Dissents:] Justice White concurred in the judgment (joined by Justices Blackmun and O'Connor, and by Justice Stevens except as to Part~I-A); Justice Blackmun concurred in the judgment; Justice Stevens concurred in the judgment (joined by Justices White and Blackmun as to Part~I). No dissents.
  \item[Chapter / Topic:] Free Speech \& Press --- Content and Viewpoint Discrimination Within Proscribable Categories; Fighting Words; Hate Speech.
  \item[Primary Source URLs:] \href{https://supreme.justia.com/cases/federal/us/505/377/}{Justia U.S. Supreme Court Center}; \href{https://www.law.cornell.edu/supremecourt/text/90-7675}{Cornell Legal Information Institute (LII)}.
\end{description}

\section{Facts and Procedural History}
In the predawn hours of June~21, 1990, petitioner R.A.V. (then a juvenile) and several other teenagers allegedly taped together broken chair legs to form a crude cross and burned it inside the fenced yard of a Black family living across the street from the home where R.A.V. was staying in St. Paul, Minnesota. Although Minnesota law provided several content-neutral criminal charges for the conduct---including terroristic threats, arson, and criminal damage to property---the City of St. Paul charged R.A.V. under its Bias-Motivated Crime Ordinance, St. Paul, Minn., Legis. Code \S~292.02 (1990). That ordinance made it a misdemeanor to place on public or private property ``a symbol, object, appellation, characterization or graffiti, including, but not limited to, a burning cross or Nazi swastika, which one knows or has reasonable grounds to know arouses anger, alarm or resentment in others on the basis of race, color, creed, religion or gender.''

R.A.V. moved to dismiss the charge, arguing that the ordinance was substantially overbroad and impermissibly content-based in violation of the First Amendment. The trial court granted the motion and dismissed the count. The Minnesota Supreme Court reversed, 464~N.W.2d~507 (Minn.~1991), rejecting the overbreadth challenge by construing the phrase ``arouses anger, alarm or resentment in others'' to reach only ``fighting words'' under \textit{Chaplinsky v. New Hampshire}, 315~U.S.~568, 572 (1942)---expression unprotected by the First Amendment. It further held that any content-based distinction survived strict scrutiny because the ordinance was narrowly tailored to serve a compelling interest in protecting the community against bias-motivated threats to public safety and order. The U.S. Supreme Court granted certiorari.

\section{Issue(s) Presented}
Whether a municipal ordinance that prohibits only a content-defined subset of ``fighting words''---those that arouse anger, alarm, or resentment on the basis of race, color, creed, religion, or gender---is facially invalid under the First Amendment.

\section{Holding}
Yes. Even accepting the Minnesota Supreme Court's narrowing construction limiting the ordinance to proscribable ``fighting words,'' the St. Paul ordinance is facially unconstitutional because it engages in impermissible content and viewpoint discrimination within a proscribable category of speech and is not necessary to serve a compelling governmental interest.

\section{Reasoning}
Writing for a five-Justice majority, Justice Scalia accepted as binding the Minnesota Supreme Court's construction of the ordinance as reaching only \textit{Chaplinsky} fighting words, and bypassed overbreadth to invalidate the law for content and viewpoint discrimination:
\begin{enumerate}
  \item \textbf{Proscribable Categories Are Not ``Invisible to the Constitution'':} Traditional categorical exceptions---obscenity, defamation, and fighting words---may be regulated because of their constitutionally proscribable content, but they are not ``entirely invisible to the Constitution.'' Like a noisy sound truck, fighting words are an unprotected \textit{mode} of communication, yet government may not regulate them based on hostility or favoritism toward the underlying nonproscribable message expressed (e.g., banning only libel critical of the government).
  \item \textbf{Three Exceptions Permitting Content Distinctions Within Proscribable Speech:} Scalia identified three narrow circumstances in which content discrimination within an unprotected category poses no realistic danger of official idea suppression:
  \begin{itemize}
    \item \textit{Reason for Proscribability:} When the basis for the content distinction consists entirely of the very reason the entire class of speech is proscribable (e.g., banning only the most patently offensive obscenity, or criminalizing only threats directed against the President under 18~U.S.C.~\S~871 because the risks of fear, disruption, and violence have special force there).
    \item \textit{Secondary Effects / Conduct Regulation:} When a content-defined subclass is associated with particular ``secondary effects'' justified without reference to expressive content (under \textit{Renton v. Playtime Theatres, Inc.}, 475~U.S.~41 (1986)), or is swept up incidentally by a general regulation directed at conduct rather than speech (such as sexually derogatory fighting words violating Title~VII's ban on employment discrimination).
    \item \textit{Catchall (No Realistic Possibility of Idea Suppression):} When the nature of the content distinction presents ``no realistic possibility that official suppression of ideas is afoot'' (e.g., banning only obscene films starring blue-eyed actresses).
  \end{itemize}
  \item \textbf{Application to St. Paul's Ordinance:} The ordinance failed this framework because it banned fighting words addressing only five disfavored topics (race, color, creed, religion, or gender) while leaving untouched equally vicious fighting words targeting political affiliation, union membership, or sexual orientation. Worse, in practical operation it produced \textit{viewpoint discrimination}: proponents of racial or religious tolerance could deploy non-enumerated fighting words freely, whereas their opponents could not invoke race- or religion-based fighting words. None of the three exceptions applied, because St. Paul targeted the emotive impact of a disfavored message rather than an especially violent mode of fighting words or content-neutral secondary effects.
  \item \textbf{Failure of Strict Scrutiny:} Although ensuring the basic human rights of historically discriminated-against groups is a compelling state interest, content discrimination was not ``reasonably necessary'' to achieve it; a content-neutral ban on \textit{all} fighting words would have produced the exact same protective effect without licensing official hostility toward particular viewpoints.
\end{enumerate}

\section{Concurring and Dissenting Opinions}
\begin{itemize}
  \item \textbf{Justice White, concurring in the judgment (joined by Justices Blackmun and O'Connor, and by Justice Stevens except Part~I-A):} Argued that the Court should have invalidated the ordinance on settled \textit{overbreadth} grounds because the Minnesota Supreme Court's construction still criminalized protected expression that merely causes ``anger, alarm or resentment.'' White sharply attacked Scalia's new ``underbreadth'' theory as abandoning the categorical approach, collapsing strict scrutiny, and inventing ad hoc exceptions that either swallowed the rule or jeopardized Title~VII hostile-workplace claims.
  \item \textbf{Justice Blackmun, concurring in the judgment:} Joined White's overbreadth analysis and warned that the majority either weakened First Amendment scrutiny across the board by dismantling categorical distinctions or issued an unprincipled aberration driven by opposition to ``politically correct speech.''
  \item \textbf{Justice Stevens, concurring in the judgment (joined by Justices White and Blackmun in Part~I):} Rejected both Scalia's near-absolute anti-content-discrimination rule and White's rigid categorical approach (Part~I-A). Favoring a contextual, multi-factor inquiry weighing content, context, and scope, Stevens argued that St. Paul's ordinance evenhandedly barred ``punches below the belt'' based on qualitatively greater social harms and would have been constitutional had it not been overbroad.
\end{itemize}

\section{Key Quotations}
\begin{quote}
``What they mean is that these areas of speech can, consistently with the First Amendment, be regulated because of their constitutionally proscribable content (obscenity, defamation, etc.)---not that they are categories of speech entirely invisible to the Constitution, so that they may be made the vehicles for content discrimination unrelated to their distinctively proscribable content.'' \\
--- Justice Scalia, \href{https://supreme.justia.com/cases/federal/us/505/377/}{505~U.S.~at~383--384} (Majority Opinion).
\end{quote}

\begin{quote}
``St. Paul has no such authority to license one side of a debate to fight freestyle, while requiring the other to follow Marquis of Queensberry rules.'' \\
--- Justice Scalia, \href{https://supreme.justia.com/cases/federal/us/505/377/}{505~U.S.~at~392} (Majority Opinion).
\end{quote}

\begin{quote}
``Let there be no mistake about our belief that burning a cross in someone's front yard is reprehensible. But St. Paul has sufficient means at its disposal to prevent such behavior without adding the First Amendment to the fire.'' \\
--- Justice Scalia, \href{https://supreme.justia.com/cases/federal/us/505/377/}{505~U.S.~at~396} (Majority Opinion).
\end{quote}

\begin{quote}
``Its decision is an arid, doctrinaire interpretation, driven by the frequently irresistible impulse of judges to tinker with the First Amendment. The decision is mischievous at best and will surely confuse the lower courts. I join the judgment, but not the folly of the opinion.'' \\
--- Justice White, \href{https://supreme.justia.com/cases/federal/us/505/377/}{505~U.S.~at~415} (Concurring in the Judgment).
\end{quote}

\section{Doctrinal Significance and Modern Legacy (Through 2026)}
\textit{R.A.V.} fundamentally reshaped modern First Amendment doctrine by establishing that the anti-viewpoint-discrimination principle penetrates even historically unprotected categories of speech, effectively foreclosing European-style ``hate speech'' bans in the United States. Its doctrinal trajectory through 2026 spans four key milestones:
\begin{itemize}
  \item \textbf{Hate-Crime Penalty Enhancements (\textit{Wisconsin v. Mitchell}):} Just one year later, a unanimous Court in \href{https://supreme.justia.com/cases/federal/us/508/476/}{\textit{Wisconsin v. Mitchell}, 508~U.S.~476 (1993)}, upheld a statute enhancing sentences for aggravated battery when the defendant intentionally selected the victim based on race. Chief Justice Rehnquist distinguished \textit{R.A.V.} on the ground that \textit{Mitchell} regulated bias-motivated \textit{conduct} (violent crime) rather than protected or proscribable \textit{expression}.
  \item \textbf{Applying \textit{R.A.V.}'s First Exception (\textit{Virginia v. Black}):} In \href{https://supreme.justia.com/cases/federal/us/538/343/}{\textit{Virginia v. Black}, 538~U.S.~343 (2003)}, the Court held that a state may ban cross burning carried out ``with an intent to intimidate'' without violating \textit{R.A.V.}, because cross burning is a particularly virulent form of ``true threat'' fitting \textit{R.A.V.}'s first exception (though a plurality struck down Virginia's provision treating cross burning as prima facie evidence of intent to intimidate).
  \item \textbf{Viewpoint Discrimination and ``Hate Speech'' (\textit{Matal v. Tam}):} In \href{https://supreme.justia.com/cases/federal/us/582/15-1293/}{\textit{Matal v. Tam}, 582~U.S.~218 (2017)}, the Court invalidated the Lanham Act's disparagement clause as unconstitutional viewpoint discrimination, reaffirming \textit{R.A.V.}'s core premise that giving offense or expressing hateful ideas cannot justify viewpoint-based government penalties.
  \item \textbf{Recent Roberts Court Refinement (2023--2026):} The Court continues to police both the boundaries of proscribable categories and viewpoint neutrality. In \textit{Counterman v. Colorado}, 600~U.S.~66 (2023), the Court required a subjective recklessness mens rea for true-threat prosecutions to preserve breathing room for speech. In \href{https://supreme.justia.com/cases/federal/us/602/22-704/}{\textit{Vidal v. Elster}, 602~U.S.~286 (2024)}, the Court distinguished historically grounded, viewpoint-neutral content rules (upholding the Lanham Act's names clause) from viewpoint discrimination. Most recently, in \href{https://supreme.justia.com/cases/federal/us/607/24-539/}{\textit{Chiles v. Salazar}, 607~U.S.~\_\_\_ (Mar.~31, 2026)}, the Court reaffirmed that only a narrow, historically pedigreed set of categories (such as fraud, defamation, and fighting words) escapes strict scrutiny, striking down viewpoint-based restrictions on professional counseling speech.
\end{itemize}

\section{Pedagogical Critiques and Exam Hypotheticals}
\textbf{Critiques:} Critics (echoing Justices White and Stevens) observe that \textit{R.A.V.} produces a striking paradox: a municipality that bans \textit{all} fighting words suppresses \textit{more} speech yet is constitutionally safer than one targeting only racial slurs, seemingly turning the First Amendment's preference for narrow tailoring upside down. Moreover, the boundary between \textit{R.A.V.}'s first exception (virulent subset of the proscribable harm) and forbidden content discrimination remains notoriously elastic, as demonstrated by the Court's later willingness in \textit{Virginia v. Black} to uphold a targeted cross-burning ban. Defenders counter that \textit{R.A.V.} prevents majorities from weaponizing broad low-value categories to selectively silence disfavored political or social viewpoints.

\textbf{Exam Hypothetical:} Following a wave of campus confrontations, a state enacts the ``Campus Harassment Act,'' making it a misdemeanor to utter face-to-face fighting words on public university grounds ``that demean a person on the basis of nationality, religion, or political ideology.'' A student prosecuted for shouting a religious slur at a classmate challenges the statute on its face under the First Amendment.
\begin{itemize}
  \item \textbf{Challenger's Argument (Under \textit{R.A.V.}):} Even if limited to \textit{Chaplinsky} fighting words, the statute selectively proscribes fighting words addressing three enumerated topics (nationality, religion, political ideology) while leaving unregulated equally provocative fighting words based on personal appearance, family, or athletics. Like St. Paul's ordinance, it fails strict scrutiny because a content-neutral ban on \textit{all} campus fighting words would equally protect campus order without singling out disfavored subjects.
  \item \textbf{Opposite Side Contrast (Distinguishing \textit{R.A.V.} via \textit{Mitchell} and \textit{Black}):} The State would argue that the statute fits \textit{R.A.V.}'s first exception (like \textit{Virginia v. Black}) if legislative findings show that nationality- and religion-based fighting words on campus carry a uniquely virulent risk of immediate violence, or that the law is evenhanded across all viewpoints (unlike St. Paul's ordinance, it covers ``political ideology'' on all sides). However, if the State instead enacted a \textit{Wisconsin v. Mitchell}--style penalty enhancement for physical battery or trespass motivated by bias---or a content-neutral ban on all true threats and fighting words---it would avoid \textit{R.A.V.} altogether.
\end{itemize}

\end{document}
